\documentclass[11pt]{article}

\usepackage[a4paper,margin=1in]{geometry}
\usepackage{amsmath,amssymb,bm}
\usepackage{booktabs}
\usepackage{hyperref}

\hypersetup{
    colorlinks=true,
    linkcolor=blue,
    citecolor=blue,
    urlcolor=blue
}

\newcommand{\eps}{\varepsilon}
\newcommand{\mean}[1]{\langle #1 \rangle}
\newcommand{\master}{geqoe\_averaged\_j2\_zonal\_note}

\title{Second-Order Near-Identity Averaging for GEqOE\\
\large Extension of the First-Order Mixed-Zonal Theory}
\author{Jose Manuel Montilla Garc\'ia}
\date{\today}

\begin{document}
\maketitle

\begin{abstract}
This companion note sketches the extension of the first-order GEqOE
mixed-zonal averaged theory (the master note, \texttt{\master{}}) to second
order in the zonal perturbation parameter. The second-order extension would
produce: (i) an $O(\eps^2)$ correction to the mean drift (including
cross-zonal products $J_n J_m$), (ii) a second-order short-period correction
$\bm u_2$ reducing the osculating reconstruction error from $O(\eps^2)$ to
$O(\eps^3)$, and (iii) a second-order inverse map further reducing the
anomaly-dependent initialization error. The central claim is that the
algebraic structure underlying the first-order direct residue
decomposition---Laurent polynomial forcing with fixed-pole-structure
Jacobian---survives at second order, making the extension computationally
demanding but not fundamentally obstructed.
\end{abstract}

\tableofcontents

\section{Motivation}

The first-order theory validated in the master note achieves sub-kilometer
position reconstruction for low-eccentricity orbits and a few kilometers for
high-eccentricity orbits ($e = 0.65$) over 10+ orbital periods. The dominant
error source is the $O(\eps^2)$ residual from neglected second-order terms
(confirmed by pointwise zonal scaling slope~$= 2.0$). A second-order theory
would push the reconstruction error to $O(\eps^3)$, which for the dominant $J_2$ perturbation
($\eps \sim 10^{-3}$) implies a reduction from hundreds of meters to
sub-meter level.

The inverse round-trip error (osculating $\to$ mean $\to$ osculating) is
currently 7.2~m (low-$e$) and 1.6~m (high-$e$). A second-order inverse map
would reduce this to centimeters, making the mean-to-osculating conversion
essentially exact for practical purposes.

\section{Second-Order Near-Identity Framework}

\subsection{Setup}

Following the notation of the master note, the exact first-order zonal
problem is
\begin{align}
\dot{\bm z} &= \eps \bm f_1(\bm z, K)
             + \eps^2 \bm f_2(\bm z, K)
             + O(\eps^3), \\
\dot K &= n_0(\bm z, K)
        + \eps g_1(\bm z, K)
        + \eps^2 g_2(\bm z, K)
        + O(\eps^3),
\end{align}
where $\bm z = (\nu, p_1, p_2, q_1, q_2)$ are the slow variables,
$\eps \sim J_n$ is the perturbation parameter, and $n_0 = w/r$ is the
unperturbed fast rate.

The first-order theory uses $\bm f_1$ and $g_1$ only. The second-order
extension requires the additional terms $\bm f_2$ and $g_2$, which arise from
expanding the exact GEqOE right-hand side to second order in the zonal
coefficients. Specifically, the linearized auxiliaries of the master note
(Section~5.3):
\begin{equation}
h = c + O(J_n),
\qquad
d = -\frac{U}{c} + O(J_n^2),
\qquad
w_h = \frac{3A}{cr^3}\,c_i\,\hat z^2 + O(J_n^2),
\end{equation}
must be carried to the next order:
\begin{align}
h &= c - \frac{r^2 U}{c} + O(J_n^2),
\label{eq:h_second}\\
d &= -\frac{U}{c} + \frac{r^2 U^2}{2c^3} + O(J_n^3),
\label{eq:d_second}\\
w_h &= w_h^{(1)} + w_h^{(2)} + O(J_n^3),
\label{eq:wh_second}
\end{align}
where $w_h^{(1)}$ is the first-order expression already used, and $w_h^{(2)}$
collects the terms quadratic in $U$ and its derivatives that arise from
expanding $h$ in $F_h$, $w_X$, and $w_Y$.

\subsection{Second-order forward map}

The near-identity transformation is extended to
\begin{align}
\bm z_{\mathrm{osc}}
&= \bar{\bm z}
+ \eps \bm u_1(\bar{\bm z}, \bar K)
+ \eps^2 \bm u_2(\bar{\bm z}, \bar K)
+ O(\eps^3),
\label{eq:forward2_z}\\
K_{\mathrm{osc}}
&= \bar K
+ \eps v_1(\bar{\bm z}, \bar K)
+ \eps^2 v_2(\bar{\bm z}, \bar K)
+ O(\eps^3),
\label{eq:forward2_K}
\end{align}
with the normalization
\begin{equation}
\mean{\bm u_2}_K = \bm 0,
\qquad
\mean{v_2}_K = 0.
\label{eq:zero_mean_2}
\end{equation}

\subsection{Second-order mean drift}

The mean drift becomes
\begin{equation}
\dot{\bar{\bm z}} = \eps \bar{\bm f}_1(\bar{\bm z})
+ \eps^2 \bar{\bm f}_2(\bar{\bm z})
+ O(\eps^3),
\end{equation}
where $\bar{\bm f}_1$ is the first-order secular drift already computed, and
$\bar{\bm f}_2$ is the second-order correction. The second-order mean drift
is determined by the solvability condition of the second-order homological
equation (see Section~\ref{sec:homological2}).

\subsection{Second-order inverse map}

The second-order inverse map is
\begin{align}
\bar{\bm z}(t_0)
&= \bm z_{\mathrm{osc}}(t_0)
- \eps \bm u_1(\bm z_{\mathrm{osc}}, K_{\mathrm{osc}})
\notag\\
&\quad
- \eps^2 \left[
\bm u_2(\bm z_{\mathrm{osc}}, K_{\mathrm{osc}})
- \frac{\partial \bm u_1}{\partial \bm z}\bigg|_{\mathrm{osc}}
  \bm u_1(\bm z_{\mathrm{osc}}, K_{\mathrm{osc}})
- \frac{\partial \bm u_1}{\partial K}\bigg|_{\mathrm{osc}}
  v_1(\bm z_{\mathrm{osc}}, K_{\mathrm{osc}})
\right]
+ O(\eps^3).
\label{eq:inverse2}
\end{align}
The additional terms at $O(\eps^2)$ correct for the error committed by
evaluating $\bm u_1$ at the osculating state instead of the mean state.
This is exactly the anomaly-dependent initialization error identified in
the master note (Section~11.6, ``Anomaly dependence''). A second-order
theory eliminates it.


\section{Second-Order Homological Equation}
\label{sec:homological2}

\subsection{Derivation}

Substituting \eqref{eq:forward2_z}--\eqref{eq:forward2_K} into the equations
of motion and collecting terms at $O(\eps^2)$ gives the second-order
homological equation for the slow variables:
\begin{equation}
\boxed{
n_0 \frac{\partial \bm u_2}{\partial \bar K}
= \bm F_2(\bar{\bm z}, \bar K)
- \bar{\bm f}_2(\bar{\bm z}),
}
\label{eq:homological2}
\end{equation}
where the effective second-order forcing is
\begin{align}
\bm F_2
&= \bm f_2(\bar{\bm z}, \bar K)
\notag\\
&\quad
+ \frac{\partial \bm f_1}{\partial \bm z}\bigg|_{\bar{\bm z}, \bar K}
  \bm u_1(\bar{\bm z}, \bar K)
+ \frac{\partial \bm f_1}{\partial K}\bigg|_{\bar{\bm z}, \bar K}
  v_1(\bar{\bm z}, \bar K)
\notag\\
&\quad
- \frac{\partial \bm u_1}{\partial \bar{\bm z}}
  \bar{\bm f}_1(\bar{\bm z})
- \frac{\partial \bm u_1}{\partial \bar K}
  \left[
    g_1(\bar{\bm z}, \bar K) - \bar g_1(\bar{\bm z})
  \right].
\label{eq:F2_forcing}
\end{align}

The solvability condition is
\begin{equation}
\bar{\bm f}_2 = \mean{\bm F_2}_K,
\label{eq:solvability2}
\end{equation}
and the solution is
\begin{equation}
\bm u_2(\bar{\bm z}, \bar K)
= \int^{\bar K}
\frac{\bar r(\phi)}{\bar w}
\left[
\bm F_2(\bar{\bm z}, \phi)
- \bar{\bm f}_2(\bar{\bm z})
\right] d\phi
- \text{(constant for zero mean)}.
\label{eq:u2_integral}
\end{equation}

\subsection{Identification of terms}

Each term in $\bm F_2$ has a clear physical origin:

\begin{enumerate}
\item $\bm f_2(\bar{\bm z}, \bar K)$: the second-order expansion of the
GEqOE right-hand side from the nonlinear auxiliaries
\eqref{eq:h_second}--\eqref{eq:wh_second}. These are the ``relaxed
linearized auxiliaries'' discussed in the motivation. They produce terms
proportional to $U^2$, $U \cdot \partial U/\partial\hat z$, etc.

\item $\frac{\partial \bm f_1}{\partial \bm z} \bm u_1$: the coupling between
the first-order forcing and the first-order short-period correction. This
is where the cross-zonal products $J_n J_m$ first enter, because
$\bm u_1$ contains contributions from all degrees $n$ while
$\bm f_1$ contains contributions from all degrees $m$.

\item $\frac{\partial \bm f_1}{\partial K} v_1$: the fast-phase feedback---the
first-order phase correction shifts the evaluation point of the forcing.

\item $-\frac{\partial \bm u_1}{\partial \bar{\bm z}} \bar{\bm f}_1$: the
secular evolution of the mean state during one orbit, correcting for the
``frozen slow state'' approximation.

\item $-\frac{\partial \bm u_1}{\partial \bar K}
[g_1 - \bar g_1]$: the interaction of the first-order short-period
correction with the perturbation to the fast rate.
\end{enumerate}


\section{Why the Algebraic Structure Survives}

\subsection{Rational structure in $F$}

The first-order theory exploits three structural facts:
\begin{enumerate}
\item[(S1)] The first-order forcing $\bm f_1$ is a finite Laurent polynomial
in $F = e^{if}$ for each $w$-harmonic.
\item[(S2)] The Jacobian $dM/dF$ has the fixed denominator
$(F+q)^2(1+qF)^2$.
\item[(S3)] The first-order short-period correction $\bm u_1$ is the
complete antiderivative of the homological equation: for zero-mean
harmonics (no secular rate), $\bm u_1$ is purely rational
(polynomial + poles at $F = -q$, $F = -1/q$, and $F = 0$).  For
secular-rate harmonics ($\Psi$\,\texttt{m=0} and $\Omega$\,\texttt{m=0}
at $n=2$; varies for $n \ge 3$), $\bm u_1$ includes a smooth
$\log(1+qF)$ correction $C_{\log}\,\varphi(K)$ where
$\varphi = \arctan(q\sin f/(1+q\cos f))$ and $C_{\log}$ is a rational
function of $(q,Q)$.  Both the rational and log parts are included in
the implementation.
\end{enumerate}

At second order, every term in $\bm F_2$ is built from:
\begin{itemize}
\item $\bm f_1$, $\bm f_2$: Laurent polynomials in $F$ (from the
trigonometric-polynomial structure of $P_n(\hat z)$ and its products).
\item $\bm u_1$, $v_1$: quasi-rational functions of $F$ with poles at $0$,
$-q$, $-1/q$ plus $\log(1+qF)$ terms for secular-rate harmonics
(from the first-order solution; see (S3) above).  If the rational-only
$\bm u_1$ is used, these are purely rational.
\item Derivatives $\partial \bm u_1/\partial \bar{\bm z}$: rational functions
of $F$ with the same pole locations (differentiation in the slow variables
does not change the $F$-pole structure).
\item Derivatives $\partial \bm f_1/\partial \bm z$: Laurent polynomials in
$F$ (differentiating a Laurent polynomial in $F$ with respect to slow
variables gives another Laurent polynomial).
\item Products of the above: rational functions of $F$ with poles only at
$0$, $-q$, $-1/q$ (products of rational functions with poles at these
locations have poles only at these locations).
\end{itemize}

\textbf{Key conclusion}: The second-order effective forcing $\bm F_2$ is a
rational function of $F$ with poles only at $F = 0$, $F = -q$, and
$F = -1/q$. The poles may be of higher order than at first order (e.g.,
triple or quadruple poles at $-q$ and $-1/q$, from the product of two
double-pole terms), but the pole \emph{locations} are unchanged.

\subsection{Consequences for the direct residue decomposition}

Since the integrand for $\bm u_2$ is
\begin{equation}
\left[\bm F_2 - \bar{\bm f}_2\right] \cdot \frac{dM}{dF}
= \frac{\text{(rational in } F\text{, poles at } 0, -q, -1/q\text{)}}
       {(F+q)^2(1+qF)^2}
\cdot (-i) \cdot \frac{(1-q^2)^3}{1+q^2} \cdot F,
\end{equation}
the full integrand is again a rational function of $F$ with poles at $0$,
$-q$, and $-1/q$. The direct residue decomposition algorithm of the master
note (Section~8) applies with these modifications:

\begin{enumerate}
\item The poles at $F = -q$ and $F = -1/q$ may now be of order 3 or 4
(instead of order 2). The residue extraction requires higher-order Laurent
coefficients, computed by repeated differentiation or by extending the
evaluation formulas of the master note.

\item The polynomial quotient from long division may be of higher degree.

\item The Taylor recurrence at $F = 0$ extends to higher shift $s$.

\item The structural mean formulas (master note, Eqs.~56--58) must be
extended to include means of higher-pole terms:
\begin{equation}
\left\langle \frac{1}{(F+q)^k} \right\rangle_M,
\qquad
\left\langle \frac{1}{(1+qF)^k} \right\rangle_M,
\qquad k = 1, 2, 3, 4.
\end{equation}
These are all computable by the same contour-integral residue calculus
(substitute $F = (z-q)/(1-qz)$ and evaluate at $z = 0$).

\item The log-term situation at second order mirrors the first-order case:
for harmonics of the second-order forcing whose mean $\bar{\bm f}_2$ is
nonzero, the mean subtraction introduces a numerator contribution at
$F = -1/q$ that does not vanish, producing logarithmic terms in the
antiderivative.  As at first order, these can be dropped in favour of
the rational-only antiderivative, which defines a self-consistent
near-identity transformation to $O(\eps^3)$ in the round-trip.
Since $\bm u_1$ now includes the complete log correction
$C_{\log}\,\varphi(K)$, the second-order cross-coupling
$\langle (\partial\bm f_1/\partial\bm z)\,\bm u_1 \rangle_M$
is exact at $O(\eps^2)$.
\end{enumerate}

\subsection{New structural mean formulas needed}

The $k$-th order mean of the pole terms can be derived systematically.
For $k = 1$ the formulas are given in the master note. For general $k$, the
substitution $F = (z-q)/(1-qz)$ gives
\begin{equation}
\frac{1}{(F+q)^k}
= \frac{(1-qz)^k}{(z(1-q^2))^k}
\cdot \frac{1}{z^0}
\cdot \ldots
\end{equation}
and the mean is determined by the coefficient of $z^0$ in the Laurent
expansion, after multiplying by the $dM/dz$ weight. The computation is
routine but produces progressively more complex rational expressions in $q$.

The means of \emph{products} of pole terms (e.g., $\langle F^j/(F+q)^k
\rangle_M$) will also be needed. These arise from the products
$\bm u_1 \cdot \partial \bm f_1/\partial \bm z$ in $\bm F_2$. Each such
mean is a single contour integral and can be evaluated in closed form, but
the total number of distinct mean formulas is larger.


\section{Cross-Zonal Products and Harmonic Content}

\subsection{Origin of cross-zonal terms}

At first order, each zonal degree $n$ contributes independently to the mean
drift and to the short-period correction. At second order, the term
$\frac{\partial \bm f_1}{\partial \bm z} \bm u_1$ creates products of
degree-$n$ forcing with degree-$m$ short-period corrections for all pairs
$(n, m)$ with $2 \le n, m \le N$.

For $N = 5$ there are 10 distinct pairs: $(2,2)$, $(2,3)$, $(2,4)$, $(2,5)$,
$(3,3)$, $(3,4)$, $(3,5)$, $(4,4)$, $(4,5)$, $(5,5)$.

\subsection{Harmonic content of the second-order mean drift}

The first-order mean drift has $\omega$-harmonic content bounded by $|m| \le
n - 2$ for each degree $n$. The cross-zonal product of degree $n$ with degree
$m$ can produce harmonics up to $|m'| \le (n-2) + (m-2) = n + m - 4$.

For $J_2 \times J_2$: only $m' = 0$ (secular). This is the classical
$J_2^2$ secular correction.

For $J_2 \times J_3$: $m' \le 1$ (odd harmonics in $\omega$). This
introduces a long-period $J_2 J_3$ term.

The total harmonic content for $N = 5$ would reach $|m'| \le 6$ in the most
extreme case ($J_5 \times J_5$, though this is the smallest in magnitude).

\subsection{The $J_2^2$ secular correction}

The most important second-order term is the $J_2^2$ secular correction to
the mean drift. This is the classical Brouwer second-order secular term. In
the present GEqOE framework, it modifies:
\begin{enumerate}
\item the precession rates $\omega_{\Psi}$ and $\omega_{\Omega}$;
\item the mean motion correction $\dot{\bar M}$;
\item and, potentially, introduces a secular drift in $\bar g$ and $\bar Q$
(which are exactly constant at first order in $J_2$).
\end{enumerate}

Whether $\bar g$ and $\bar Q$ acquire a $J_2^2$ secular drift depends on
whether the $J_2^2$ averaged disturbing function retains $\omega$-dependence
through the coupling terms. This is a key question that the symbolic
computation must answer.


\section{Computational Feasibility}

\subsection{Expression complexity}

The first-order theory required $\sim$80 minutes of one-time symbolic
computation (dominated by $J_5$ Psi and Omega). At second order, the
complexity grows in three ways:

\begin{enumerate}
\item \textbf{Number of terms}: 10 cross-zonal pairs instead of 4 isolated
degrees, and each pair involves products of the (already large) first-order
expressions.

\item \textbf{Polynomial degree}: The products $\bm u_1 \cdot
\partial \bm f_1/\partial \bm z$ have roughly double the Laurent polynomial
support of the first-order forcing. The polynomial long division and
Taylor recurrence steps scale polynomially with this degree.

\item \textbf{Coefficient complexity}: The rational coefficients in
$\mathbb{Q}(q, Q)$ grow in total degree. GCD simplification becomes
progressively more expensive.
\end{enumerate}

\textbf{Rough estimate}: The dominant cost is the symbolic simplification of
rational expressions in $(q, Q)$. If the first-order $J_5$ computation took
$\sim$80 minutes, the second-order computation might take days to weeks of
CPU time. This is large but still one-time and parallelizable (each
cross-zonal pair and each slow variable can be computed independently).

\subsection{Mitigation strategies}

\begin{enumerate}
\item \textbf{Coefficient-wise operations}: The coefficient-wise scaling
optimization from the first-order theory (master note, Section~8.4) would be
even more critical at second order. All multiplications, divisions, and GCD
computations should be performed coefficient-by-coefficient in the $F$-Laurent
expansion, never on the combined rational expression.

\item \textbf{Modular arithmetic}: For verification, the symbolic expressions
could be evaluated modulo large primes before full rational reconstruction.
This would catch errors early without the cost of full symbolic
simplification.

\item \textbf{Incremental validation}: Each cross-zonal pair $(n, m)$ can be
validated independently against numerical second-order differences before
combining into the full second-order model.

\item \textbf{Truncation by magnitude}: The $J_5 \times J_5$ and
$J_4 \times J_5$ cross terms are negligibly small. In practice, only the
$(2,2)$, $(2,3)$, $(2,4)$, $(3,3)$ pairs are likely to matter, reducing the
computation to 4--6 pairs.
\end{enumerate}

\subsection{Validation strategy}

The second-order theory can be validated incrementally:

\begin{enumerate}
\item \textbf{$J_2^2$ only}: Implement the $(2,2)$ second-order correction
and validate against the osculating Taylor integrator with $J_2$ only. The
reconstruction error should drop from $O(J_2)$ to $O(J_2^2)$, i.e., from
$\sim$200~m to $\sim$sub-meter for the low-$e$ test case.

\item \textbf{Zonal scaling}: The scaling test should now show a log-log
slope of $\approx 2$ (instead of $\approx 1$), confirming that the dominant
residual is $O(\eps^2)$ as expected.

\item \textbf{Cross-zonal pairs}: Add $(2,3)$, $(2,4)$, etc., one at a time
and verify that the reconstruction error decreases (or at least does not
increase) with each addition.

\item \textbf{Second-order inverse map}: The round-trip error
(osc $\to$ mean $\to$ osc) should drop from meters to centimeters.
\end{enumerate}


\section{Expected Accuracy Improvement}

For the low-eccentricity test case ($e = 0.05$, $a = 9000$~km):

\begin{center}
\begin{tabular}{lcc}
\toprule
Metric & First-order & Second-order (projected) \\
\midrule
Position RMS & $\sim$200 m & $\sim$0.1--1 m \\
Round-trip error & 7.2 m & $\sim$cm \\
Scaling slope (pointwise) & 2.0 & $\approx 3$ \\
\bottomrule
\end{tabular}
\end{center}

For the high-eccentricity case ($e = 0.65$), the improvement factor is
similar but the absolute errors are larger due to the pole-proximity effect.
The $O(\eps^2)$ residual would still be tens of meters rather than sub-meter,
because the first-order short-period amplitudes are already large ($q \approx
0.37$).


\section{Outline of Implementation Steps}

\begin{enumerate}
\item Derive the second-order expansions \eqref{eq:h_second}--\eqref{eq:wh_second}
of the GEqOE auxiliaries. This gives $\bm f_2$ and $g_2$.

\item Compute the derivatives $\partial \bm f_1/\partial \bm z$ and
$\partial \bm f_1/\partial K$ symbolically in the $(g, Q, \Psi, \Omega, G)$
parameterization.

\item For each cross-zonal pair $(n, m)$, compute the products
$(\partial \bm f_1^{(n)}/\partial \bm z) \cdot \bm u_1^{(m)}$ and expand in
$w$-harmonics and $F$-Laurent series.

\item Apply the direct residue decomposition (with extended pole orders) to
each resulting integrand.

\item Compute the second-order mean drift $\bar{\bm f}_2$ via the
solvability condition \eqref{eq:solvability2}.

\item Assemble the second-order short-period correction $\bm u_2$ and the
second-order inverse map \eqref{eq:inverse2}.

\item Validate incrementally as described in Section~6.3.
\end{enumerate}


\section{Conclusion}

The second-order extension of the GEqOE averaged theory is algebraically
tractable because the rational-in-$F$ structure with fixed pole locations
survives at second order. The main challenges are computational (expression
complexity, symbolic GCD cost) rather than structural. The expected payoff is
a reduction of the reconstruction error from $O(\eps)$ to $O(\eps^2)$,
bringing the analytical theory into sub-meter accuracy for low-eccentricity
orbits. The implementation is parallelizable and can be validated
incrementally, starting with the most important $J_2^2$ secular correction.

\end{document}
